% ============================================================
%  Chapter 3 -- Exercises 3.1-3.25
% ============================================================
\rsection{Exercises}

\begin{roadmap}[How to use these]
Chapter 3's exercise set contains two of the most important constructions in
the book. \emph{3.23--3.24} build the \textbf{completion} of an arbitrary
metric space out of Cauchy sequences --- Tao's construction of $\R$ itself,
carried out here as an exercise, with 3.25 closing the circle by asking what
the completion of $\Q$ is. \emph{3.22} is \textbf{Baire's theorem} in its
general form, promised at Exercise 2.30.

\medskip
\raggedright
\textbf{Technique builders.} \emph{3.3} and \emph{3.16--3.18} are the
monotone-convergence workhorses: recursive sequences squeezed by 3.14, with
3.16 being Newton's method for square roots and its blazing quadratic
convergence. \emph{3.14} develops Ces\`aro means --- the first taste of
summability theory, needed for Fourier series (8.15--8.16) and Tauberian
arguments.

\smallskip
\textbf{Counterexample duty.} \emph{3.6--3.8} and \emph{3.11--3.12} sharpen
the convergence tests; 3.11(b) and 3.12(b) are the pairs Rudin cited in the
text to kill the ``boundary between convergence and divergence'' conjecture.

\smallskip
\textbf{Routine.} 3.1, 3.2, 3.9, 3.20 are quick checks; do them first.
\end{roadmap}

\begin{enumerate}[label=\textbf{3.\arabic*.}, leftmargin=3.0em, itemsep=9pt,
                  topsep=8pt, labelsep=0.5em]

\item Prove that convergence of $(s_n)$ implies convergence of $(|s_n|)$. Is
the converse true?\kn{Reverse triangle inequality (Exercise 1.13):
$\bigl||s_n|-|s|\bigr| \le |s_n-s|$. The converse fails at once:
$s_n = (-1)^n$.}

\item Calculate $\lim_{n\to\infty} (\sqrt{n^2+n}-n)$.\kn{Multiply and divide
by the conjugate $\sqrt{n^2+n}+n$; the answer is $\tfrac12$. The moral:
$\infty - \infty$ expressions are settled by algebra first, limit theorems
second --- 3.3 applies only after the difference is rewritten as a quotient.}

\item If $s_1 = \sqrt2$ and
\[ s_{n+1} = \sqrt{2+\sqrt{s_n}} \qquad (n = 1,2,3,\dots), \]
prove that $(s_n)$ converges, and that $s_n < 2$ for
$n = 1,2,3,\dots$.\kn{The 3.14 template, used again in 3.16--3.18: prove
boundedness ($s_n < 2$) and monotonicity together by induction, then monotone
convergence hands you a limit. Only afterwards may you pass to the limit in
the recursion to identify it --- existence must precede the algebra.}

\item Find the upper and lower limits of the sequence $(s_n)$ defined by
\[ s_1 = 0; \qquad s_{2m} = \frac{s_{2m-1}}{2}; \qquad
   s_{2m+1} = \frac12 + s_{2m}. \]\kn{Track the two subsequences separately:
the even-indexed terms crawl up toward $\tfrac12$, the odd ones toward $1$.
So $E = \{\tfrac12, 1\}$, $\liminf = \tfrac12$, $\limsup = 1$ --- a sequence
need not oscillate symmetrically about anything.}

\item For any two real sequences $(a_n)$, $(b_n)$, prove that
\[ \limsup_{n\to\infty} (a_n+b_n) \le
   \limsup_{n\to\infty} a_n + \limsup_{n\to\infty} b_n, \]
provided the sum on the right is not of the form
$\infty-\infty$.\kw{Inequality, not equality --- and strictness is common:
$a_n = (-1)^n$, $b_n = (-1)^{n+1}$ gives $0 < 2$. The two sequences need not
peak at the same indices. (The excluded case is exactly the arithmetic left
undefined at 1.23.) The working proof uses the tail-sup form from the box at
3.17: $\sup_{n\ge N}(a_n+b_n) \le \sup_{n\ge N}a_n + \sup_{n\ge N}b_n$, then
let $N\to\infty$.}

\item Investigate the behavior (convergence or divergence) of $\Sigma a_n$ if
\begin{enumerate}[label=(\alph*), leftmargin=2.2em, itemsep=2pt, topsep=3pt]
  \item $a_n = \sqrt{n+1}-\sqrt n$;
  \item $a_n = \dfrac{\sqrt{n+1}-\sqrt n}{n}$;
  \item $a_n = (\sqrt[n]{n}-1)^n$;
  \item $a_n = \dfrac{1}{1+z^n}$, for complex values of $z$.
\end{enumerate}
\kn{(a) telescopes to $\sqrt{N+1}-1$: diverges. (b) conjugate trick gives
$a_n \sim \tfrac12 n^{-3/2}$: converges by comparison with 3.28. (c) root
test: $\sqrt[n]{a_n} = \sqrt[n]{n}-1 \to 0$ by 3.20(c): converges. (d) splits
by $|z|$: for $|z|>1$ comparison with $\Sigma|z|^{-n}$ converges; for
$|z| \le 1$ the terms do not tend to $0$ (3.23) --- watch $|z|=1$, where
$1+z^n$ can even vanish.}

\item Prove that the convergence of $\Sigma a_n$ implies the convergence of
\[ \sum_{n=1}^{\infty} \frac{\sqrt{a_n}}{n}, \]
if $a_n \ge 0$.\kn{The trick to remember:
$\sqrt{a_n}/n \le \tfrac12(a_n + 1/n^2)$ by the AM--GM inequality
$xy \le \tfrac12(x^2+y^2)$ --- which is Schwarz 1.35 for $n=1$. Both halves
converge, so comparison finishes it.}

\item If $\Sigma a_n$ converges and if $(b_n)$ is monotonic and bounded,
prove that $\Sigma a_nb_n$ converges.\kn{Summation by parts, but not quite
3.42, whose hypothesis $b_n \to 0$ may fail here. Write $b_n = (b_n-b)+b$
where $b = \lim b_n$ (exists by 3.14): the $b$-part converges by 3.47, and
$(b_n-b)$ decreases (or increases) to $0$, so 3.42 handles the rest. Note the
partial sums $A_n$ are bounded because $\Sigma a_n$ \emph{converges} --- a
weaker demand than 3.42 makes.}

\item Find the radius of convergence of each of the following power series:
\begin{multicols}{2}
\begin{enumerate}[label=(\alph*), leftmargin=2.2em, itemsep=2pt, topsep=3pt]
  \item $\displaystyle\sum n^3z^n$,
  \item $\displaystyle\sum \frac{2^n}{n!}z^n$,
  \item $\displaystyle\sum \frac{2^n}{n^2}z^n$,
  \item $\displaystyle\sum \frac{n^3}{3^n}z^n$.
\end{enumerate}
\end{multicols}
\kn{All four by 3.37's byproduct ($\lim c_{n+1}/c_n$, when it exists, is
$\alpha$): (a) $R=1$ --- polynomial factors never move $R$, by 3.20(c);
(b) $R=\infty$; (c) $R=\tfrac12$; (d) $R=3$.}

\item Suppose that the coefficients of the power series $\Sigma a_nz^n$ are
integers, infinitely many of which are distinct from zero. Prove that the
radius of convergence is at most $1$.\kn{A nonzero integer has
$|a_n| \ge 1$, so $\sqrt[n]{|a_n|} \ge 1$ infinitely often and
$\alpha = \limsup \ge 1$. Integrality forces smallness to be impossible ---
the same arithmetic leverage as in the irrationality proof 3.32.}

\item Suppose $a_n > 0$, $s_n = a_1+\cdots+a_n$, and $\Sigma a_n$ diverges.
\begin{enumerate}[label=(\alph*), leftmargin=2.2em, itemsep=3pt, topsep=3pt]
  \item Prove that $\displaystyle\sum \frac{a_n}{1+a_n}$ diverges.
  \item Prove that
        \[ \frac{a_{N+1}}{s_{N+1}} + \cdots + \frac{a_{N+k}}{s_{N+k}}
           \ge 1 - \frac{s_N}{s_{N+k}} \]
        and deduce that $\displaystyle\sum \frac{a_n}{s_n}$ diverges.
  \item Prove that
        \[ \frac{a_n}{s_n^2} \le \frac{1}{s_{n-1}} - \frac{1}{s_n} \]
        and deduce that $\displaystyle\sum \frac{a_n}{s_n^2}$ converges.
  \item What can be said about
        \[ \sum \frac{a_n}{1+na_n} \qquad\text{and}\qquad
           \sum \frac{a_n}{1+n^2a_n}\,? \]
\end{enumerate}
\kd{This is the divergent half of Rudin's ``no boundary'' claim after 3.29:
by (b), $\Sigma a_n/s_n$ diverges yet its terms are smaller than $a_n$ by the
unbounded factor $s_n$ --- every divergent series dominates a slower one. Part
(b)'s display is a Cauchy-criterion violation served ready-made: fix $N$, let
$k$ grow, and the blocks refuse to shrink because $s_{N+k}\to\infty$. Part
(c) telescopes. In (d) the first can do either (try $a_n=1$ vs.\ $a_n$
supported on sparse blocks); the second always converges, since
$a_n/(1+n^2a_n) \le 1/n^2$.}

\item Suppose $a_n > 0$ and $\Sigma a_n$ converges. Put
\[ r_n = \sum_{m=n}^{\infty} a_m. \]
\begin{enumerate}[label=(\alph*), leftmargin=2.2em, itemsep=3pt, topsep=3pt]
  \item Prove that
        \[ \frac{a_m}{r_m} + \cdots + \frac{a_n}{r_n} > 1 - \frac{r_n}{r_m} \]
        if $m < n$, and deduce that $\displaystyle\sum \frac{a_n}{r_n}$
        diverges.
  \item Prove that
        \[ \frac{a_n}{\sqrt{r_n}} < 2(\sqrt{r_n}-\sqrt{r_{n+1}}) \]
        and deduce that $\displaystyle\sum \frac{a_n}{\sqrt{r_n}}$
        converges.
\end{enumerate}
\kd{The convergent half of ``no boundary'': by (b), dividing a convergent
series by the vanishing factor $\sqrt{r_n}$ leaves it convergent --- every
convergent series is dominated by a slower one. Together with 3.11(b) this
kills the conjecture in both directions. The engine in both parts is the same:
$r_n$ telescopes ($a_n = r_n - r_{n+1}$), and (b) is the algebraic identity
$\frac{r_n-r_{n+1}}{\sqrt{r_n}} < 2(\sqrt{r_n}-\sqrt{r_{n+1}})$, i.e.\
$\sqrt{r_n}+\sqrt{r_{n+1}} < 2\sqrt{r_n}$.}

\item Prove that the Cauchy product of two absolutely convergent series
converges absolutely.\kn{Compare $\sum_n|c_n|$ against
$(\Sigma|a_n|)(\Sigma|b_n|)$: every partial sum of $|c_n|$ is trapped in a
triangle of the array of $|a_kb_j|$, which sits inside a full rectangle. With
3.50 this closes the multiplication story: absolute $\times$ absolute is
absolute; absolute $\times$ conditional converges to the right value
(Mertens); conditional $\times$ conditional can diverge (3.49).}

\item If $(s_n)$ is a complex sequence, define its arithmetic means
$\sigma_n$ by
\[ \sigma_n = \frac{s_0 + s_1 + \cdots + s_n}{n+1} \qquad (n = 0,1,2,\dots). \]
\begin{enumerate}[label=(\alph*), leftmargin=2.2em, itemsep=3pt, topsep=3pt]
  \item If $\lim s_n = s$, prove that $\lim \sigma_n = s$.
  \item Construct a sequence $(s_n)$ which does not converge, although
        $\lim \sigma_n = 0$.
  \item Can it happen that $s_n > 0$ for all $n$ and that
        $\limsup s_n = \infty$, although $\lim \sigma_n = 0$?
  \item Put $a_n = s_n - s_{n-1}$ for $n \ge 1$. Show that
        \[ s_n - \sigma_n = \frac{1}{n+1}\sum_{k=1}^{n} ka_k. \]
        Assume that $\lim (na_n) = 0$ and that $(\sigma_n)$ converges. Prove
        that $(s_n)$ converges. [This gives a converse of (a), but under the
        additional assumption that $na_n \to 0$.]
  \item Derive the last conclusion from a weaker hypothesis: Assume
        $M < \infty$, $|na_n| \le M$ for all $n$, and
        $\lim \sigma_n = \sigma$. Prove that $\lim s_n = \sigma$, by
        completing the following outline:

        If $m < n$, then
        \[ s_n - \sigma_n = \frac{m+1}{n-m}(\sigma_n - \sigma_m)
           + \frac{1}{n-m}\sum_{i=m+1}^{n} (s_n - s_i). \]
        For these $i$,
        \[ |s_n - s_i| \le \frac{(n-i)M}{i+1} \le \frac{(n-m-1)M}{m+2}. \]
        Fix $\varepsilon > 0$ and associate with each $n$ the integer $m$ that
        satisfies
        \[ m \le \frac{n-\varepsilon}{1+\varepsilon} < m+1. \]
        Then $(m+1)/(n-m) \le 1/\varepsilon$ and $|s_n - s_i| < M\varepsilon$.
        Hence
        \[ \limsup_{n\to\infty} |s_n - \sigma| \le M\varepsilon. \]
        Since $\varepsilon$ was arbitrary, $\lim s_n = \sigma$.
\end{enumerate}
\kd{Ces\`aro summability, and a first look at an entire subject. Part (a) says
averaging never destroys a limit; (b), with $s_n = (-1)^n$, says it can
manufacture one --- so ``Ces\`aro sum'' genuinely extends ``sum,'' which is
how Fourier series will be rescued at 8.15--8.16 (Fej\'er's idea). Parts (d)
and (e) are \emph{Tauberian} theorems: conditions under which the extension
collapses back to ordinary convergence. Hardy's condition $|na_n| \le M$ in
(e) is the historical original. The two-parameter bookkeeping in the outline
--- $m$ tied to $n$ through $\varepsilon$ --- is 3.31's freeze-and-release
discipline pushed one level up.}

\item Definition 3.21 can be extended to the case in which the $a_n$ lie in
some fixed $\R^k$. Absolute convergence is defined as convergence of
$\Sigma|\mathbf{a}_n|$. Show that Theorems 3.22, 3.23, 3.25(a), 3.33, 3.34,
3.42, 3.45, 3.47, and 3.55 are true in this more general setting. (Only
slight modifications are required in any of the proofs.)\kn{The checklist of
which theory is metric and which is order-theoretic: everything driven by the
Cauchy criterion and the triangle inequality survives in $\R^k$ verbatim. What
is missing from the list matters too --- 3.24, 3.27, and comparison (b) need
\emph{nonnegativity}, and the root/ratio tests survive only because they test
the scalar series $\Sigma|\mathbf{a}_n|$.}

\item Fix a positive number $\alpha$. Choose $x_1 > \sqrt\alpha$, and define
$x_2, x_3, x_4, \dots$ by the recursion formula
\[ x_{n+1} = \frac12\left(x_n + \frac{\alpha}{x_n}\right). \]
\begin{enumerate}[label=(\alph*), leftmargin=2.2em, itemsep=3pt, topsep=3pt]
  \item Prove that $(x_n)$ decreases monotonically and that
        $\lim x_n = \sqrt\alpha$.
  \item Put $\varepsilon_n = x_n - \sqrt\alpha$, and show that
        \[ \varepsilon_{n+1} = \frac{\varepsilon_n^2}{2x_n}
           < \frac{\varepsilon_n^2}{2\sqrt\alpha}, \]
        so that, setting $\beta = 2\sqrt\alpha$,
        \[ \varepsilon_{n+1} < \beta\left(\frac{\varepsilon_1}{\beta}\right)^{2^n}
           \qquad (n = 1,2,3,\dots). \]
  \item This is a good algorithm for computing square roots, since the
        recursion formula is simple and the convergence is extremely rapid.
        For example, if $\alpha = 3$ and $x_1 = 2$, show that
        $\varepsilon_1/\beta < 1/10$ and that therefore
        \[ \varepsilon_5 < 4\cdot10^{-16}, \qquad
           \varepsilon_6 < 4\cdot10^{-32}. \]
\end{enumerate}
\kn{Newton's method on $f(x) = x^2 - \alpha$, and the recursion Chapter 1's
note on Example 1.1 promised: it always lands \emph{above} $\sqrt\alpha$
(AM--GM), so the sequence is decreasing and 3.14 applies. Part (b) is
quadratic convergence --- the error squares each step, doubling the correct
digits --- against which the linear rate of 1.1's map $q = (2p+2)/(p+2)$, and
of 3.17 below, can be measured.}

\item Fix $\alpha > 1$. Take $x_1 > \sqrt\alpha$, and define
\[ x_{n+1} = \frac{\alpha+x_n}{1+x_n} = x_n + \frac{\alpha-x_n^2}{1+x_n}. \]
\begin{enumerate}[label=(\alph*), leftmargin=2.2em, itemsep=2pt, topsep=3pt]
  \item Prove that $x_1 > x_3 > x_5 > \cdots$.
  \item Prove that $x_2 < x_4 < x_6 < \cdots$.
  \item Prove that $\lim x_n = \sqrt\alpha$.
  \item Compare the rapidity of convergence of this process with the one
        described in Exercise 3.16.
\end{enumerate}
\kn{The general-$\alpha$ version of Chapter 1's map (there $\alpha = 2$:
$q = (2+p)/(1+p)$). Unlike Newton, it alternates sides of $\sqrt\alpha$, so
the two subsequences must be handled separately by 3.14 and then reunited ---
the strategy of 3.4. Convergence is linear: the error is multiplied by roughly
$(\sqrt\alpha-1)/(\sqrt\alpha+1)$ each step, one new digit per few
iterations, versus Newton's doubling.}

\item Replace the recursion formula of Exercise 3.16 by
\[ x_{n+1} = \frac{p-1}{p}x_n + \frac{\alpha}{p}x_n^{-p+1}, \]
where $p$ is a fixed positive integer, and describe the behavior of the
resulting sequences $(x_n)$.\kn{Newton for $f(x) = x^p - \alpha$: the limit is
$\sqrt[p]{\alpha}$, with the same quadratic convergence and the same
monotone-decrease-from-above proof. A numerical companion to Theorem 1.21,
which proved the $p$th root exists; this computes it.}

\item Associate to each sequence $a = (\alpha_n)$, in which $\alpha_n$ is $0$
or $2$, the real number
\[ x(a) = \sum_{n=1}^{\infty} \frac{\alpha_n}{3^n}. \]
Prove that the set of all $x(a)$ is precisely the Cantor set described in
Theorem 2.44.\kn{The address system of Chapter 2's binary-tree figure made
exact: digit $\alpha_n$ chooses the left or right third at stage $n$, and the
series is the point the nested thirds close on. Series notation turns the
geometric description into a formula --- and makes the bijection with
$\{0,2\}^\N \sim \{0,1\}^\N$ (hence uncountability, 2.14) immediate.}

\item Suppose $(p_n)$ is a Cauchy sequence in a metric space $X$, and some
subsequence $(p_{n_i})$ converges to a point $p \in X$. Prove that the full
sequence $(p_n)$ converges to $p$.\kn{Split $d(p_n,p) \le d(p_n,p_{n_i}) +
d(p_{n_i},p)$: the first half is small because the sequence is Cauchy, the
second because the subsequence converges. Small but used constantly ---
including in the very next two exercises, and it is how 3.11(b) could have
been proved from 3.6(a) directly.}

\item Prove the following analogue of Theorem 3.10(b): If $(E_n)$ is a
sequence of closed, nonempty, and bounded sets in a \emph{complete} metric
space $X$, if $E_n \supset E_{n+1}$, and if
\[ \lim_{n\to\infty} \operatorname{diam} E_n = 0, \]
then $\bigcap_{n=1}^\infty E_n$ consists of exactly one point.\kn{Compactness
weakened to closed-and-bounded, compensated by completeness: pick
$x_n \in E_n$; shrinking diameters make $(x_n)$ Cauchy, completeness gives
the limit, closedness puts it in every $E_n$. Compare 2.36, where compactness
did the work with no diameter hypothesis at all --- and note that here,
unlike there, ``bounded'' can even be dropped once diameters shrink.}

\item Suppose $X$ is a nonempty complete metric space, and $(G_n)$ is a
sequence of dense open subsets of $X$. Prove Baire's theorem, namely, that
$\bigcap_{n=1}^\infty G_n$ is not empty. (In fact, it is dense in $X$.)
\emph{Hint:} Find a shrinking sequence of neighborhoods $E_n$ such that
$\overline{E}_n \subset G_n$, and apply Exercise 3.21.\kd{The general Baire
theorem, upgrading Exercise 2.30 from $\R^k$ to any complete metric space ---
with 3.21 replacing compactness as the engine. The construction: density of
$G_1$ gives a ball inside it; density of $G_2$ gives a smaller ball inside
both, radius halved; iterate, and 3.21 yields a point surviving every stage.
This theorem underlies uniform boundedness and the open mapping theorem in
functional analysis, and Exercise 2.30's consequences (e.g.\ $\R$
uncountable) now hold in every complete metric space without isolated
points.}

\item Suppose $(p_n)$ and $(q_n)$ are Cauchy sequences in a metric space $X$.
Show that the sequence $(d(p_n,q_n))$ converges. \emph{Hint:} For any $m$,
$n$,
\[ d(p_n,q_n) \le d(p_n,p_m) + d(p_m,q_m) + d(q_m,q_n); \]
it follows that
\[ |d(p_n,q_n) - d(p_m,q_m)| \]
is small if $m$ and $n$ are large.\kn{Note where the limit lives: the
\emph{real} sequence $(d(p_n,q_n))$ is Cauchy, so it converges by 3.11(c) ---
completeness of $\R$ --- even though $(p_n)$ and $(q_n)$ may converge to
nothing in $X$. That asymmetry is exactly what the next exercise exploits.}

\item Let $X$ be a metric space.
\begin{enumerate}[label=(\alph*), leftmargin=2.2em, itemsep=3pt, topsep=3pt]
  \item Call two Cauchy sequences $(p_n)$, $(q_n)$ in $X$ \emph{equivalent} if
        \[ \lim_{n\to\infty} d(p_n,q_n) = 0. \]
        Prove that this is an equivalence relation.
  \item Let $X^*$ be the set of all equivalence classes so obtained. If
        $P \in X^*$, $Q \in X^*$, $(p_n) \in P$, $(q_n) \in Q$, define
        \[ \Delta(P,Q) = \lim_{n\to\infty} d(p_n,q_n); \]
        by Exercise 3.23, this limit exists. Show that the number
        $\Delta(P,Q)$ is unchanged if $(p_n)$ and $(q_n)$ are replaced by
        equivalent sequences, and hence that $\Delta$ is a distance function
        in $X^*$.
  \item Prove that the resulting metric space $X^*$ is complete.
  \item For each $p \in X$, there is a Cauchy sequence all of whose terms are
        $p$; let $P_p$ be the element of $X^*$ which contains this sequence.
        Prove that
        \[ \Delta(P_p,P_q) = d(p,q) \]
        for all $p,q \in X$. In other words, the mapping $\varphi$ defined by
        $\varphi(p) = P_p$ is an \emph{isometry} (i.e., a distance-preserving
        mapping) of $X$ into $X^*$.
  \item Prove that $\varphi(X)$ is dense in $X^*$, and that
        $\varphi(X) = X^*$ if $X$ is complete. By (d), we may identify $X$
        and $\varphi(X)$ and thus regard $X$ as embedded in the complete
        metric space $X^*$. We call $X^*$ the \emph{completion} of $X$.
\end{enumerate}
\kd{The construction the appendix to Chapter 1 declined: a missing limit is
\emph{represented by the sequences that reach for it}, with two sequences
identified when they reach for the same thing. This is Cantor's route to $\R$
and the one Tao's \emph{Analysis I} takes as its foundation (his Chapter 5:
reals \emph{are} equivalence classes of rational Cauchy sequences). Note what
makes it outrun Dedekind cuts, as the appendix's note anticipated: nothing
here mentions order --- only $d$ --- so \emph{every} metric space acquires a
completion, function spaces included. Part (c) has the construction's one
subtlety: a Cauchy sequence \emph{of equivalence classes} is unwound by
picking representatives and a diagonal argument, the 3.7-style chase again.}

\item Let $X$ be the metric space whose points are the rational numbers, with
the metric $d(x,y) = |x-y|$. What is the completion of this space? (Compare
Exercise 3.24.)\kn{$\R$, of course --- and the point of asking is that the
completion was built \emph{without} $\R$: 3.24 used only $\Q$'s own metric.
Chapter 1 characterised $\R$ by order (least upper bounds); this exercise
constructs it from distance. Same object, two foundations, and the
uniqueness-up-to-isomorphism noted in Step 9 of the appendix is what
guarantees they agree.}

\end{enumerate}
